\begin{table}[!htbp]\small\setlength{\tabcolsep}{4pt}
\caption{Recettes des impôts concernés et solde public (Md€ de 2025)}\label{tab:t1}
\begin{tabularx}{\linewidth}{@{}Xrrrr@{}}
\toprule
\textbf{Mesure et impôts concernés} & \textbf{Recettes 2025} & \textbf{Après réforme} & \textbf{Solde 2032} & \textbf{Solde à terme} \\
\midrule
\textbf{1. Nouvel impôt sur le revenu} & \textbf{260,2} & \textbf{260,2} & \textbf{0,0} & \textbf{0,0} \\
\quad Impôt sur le revenu, net des crédits d'impôt & 94,3 & fusionnés & & \\
\quad CSG & 156,6 & & & \\
\quad CRDS & 9,3 & & & \\
\textbf{2. Fiscalité foncière} & \textbf{61,8} & \textbf{67,3} & \textbf{+5,2} & \textbf{+5,4} \\
\quad Taxe foncière communale (sol et bâti) & 44,2 & 44,2 & 0,0 & 0,0 \\
\quad Taxe nationale sur la valeur du sol & 0,0 & 15,0 & +15,0 & +15,0 \\
\quad Droits de mutation à titre onéreux & 17,7 & 8,1 & $-$8,1 & $-$8,1 \\
\quad Déductibilité de la taxe sur le sol, calendrier & & & $-$1,8 & $-$1,5 \\
\textbf{3. Impôt sur les sociétés} & \textbf{59,3} & \textbf{59,2} & \textbf{$-$5,7} & \textbf{$-$0,1} \\
\textbf{4. Prélèvement forfaitaire unique (inchangé)} & & & \textbf{0,0} & \textbf{0,0} \\
\midrule \textbf{Total} & \textbf{381,3} & \textbf{386,6} & \textbf{$-$0,6} & \textbf{+5,3} \\
\bottomrule
\end{tabularx}
\par\vspace{2pt}{\raggedright\scriptsize\color{gris}Source : Insee, comptes nationaux 2025, tableau 3.217 (recettes 2025) ; calculs des auteurs. Après réforme et solde à terme : régime de croisière, une fois achevée la transition vers la déduction immédiate des investissements (recettes de 2040 rapportées à l'économie de 2025). Les droits de mutation de 2025 comprennent la hausse temporaire autorisée jusqu'en 2028 ; le solde est calculé par rapport au taux permanent (16,1 Md€). L'effet sur les prélèvements obligatoires est égal au solde public.\par}
\end{table}
